% 这个文件不要修改-- do not modify this file please.


\usepackage[utf8]{inputenc}
\usepackage[T1]{fontenc} % 建议加上，提升字体编码支持

\usepackage[margin=1in]{geometry}
\usepackage[most]{tcolorbox}
\usepackage{graphicx}
\usepackage{float}
\usepackage{enumitem}

% --- 数学环境 ---
\usepackage{amsmath}
\usepackage{amssymb}
\usepackage{amsthm}
\usepackage{thmtools}

% --- 算法 ---
\usepackage{algorithm}
\usepackage{algpseudocode}

% --- 引用与链接 (顺序很关键) ---
\usepackage{hyperref}

\hypersetup{
    colorlinks=true,
    linkcolor={red!50!black},  % 内部链接（图表、章节）：深红色
    citecolor={blue!50!black}, % 参考文献：深蓝色
    urlcolor={blue!80!black},  % 外部链接：偏亮的深蓝
    breaklinks=true            % 允许超链接折行
}
\usepackage[capitalize,nameinlink]{cleveref}

\usepackage{xspace}
\usepackage{xcolor}
\usepackage{blindtext}
\usepackage{url}
\usepackage{verbatim}

\usepackage[style=alphabetic, backend=biber, maxcitenames=99, maxbibnames=99, minalphanames=3, labelalpha=true]{biblatex}

\theoremstyle{plain}
\newtheorem{theorem}{Theorem}[section]
\newtheorem{lemma}[theorem]{Lemma}
\newtheorem{corollary}[theorem]{Corollary}
\newtheorem{conjecture}[theorem]{Conjecture}
\newtheorem{claim}[theorem]{Claim}
\theoremstyle{definition}
\newtheorem{definition}[theorem]{Definition}
\newtheorem{example}[theorem]{Example}
\newtheorem{question}[theorem]{Question}
\newtheorem{problem}[theorem]{Problem}
\newtheorem{open}[theorem]{Open problem}
\theoremstyle{remark}
\newtheorem{remark}[theorem]{Remark}

\newcommand{\PPAD}{\ensuremath{\mathsf{PPAD}}\xspace}
\newcommand{\TFNP}{\ensuremath{\mathsf{TFNP}}\xspace}
\newcommand{\NP}{\ensuremath{\mathsf{NP}}\xspace}
\newcommand{\cmt}[1]{} % 保持不变，确认你确实想让注释隐藏
